\def\writeANS{}
\begin{loigiaibth}{47}

\end{loigiaibth}
\def\writeANS{}
\begin{loigiaibth}{48}
 Bao gồm 2, 3.
\end{loigiaibth}
\def\writeANS{}
\begin{loigiaibth}{49}
 Bao gồm 2, 3, 4, 5.
\end{loigiaibth}
\def\writeANS{}
\begin{loigiaibth}{50}
  $n_{H_2} = 1{,}5$ mol $\Rightarrow $ Q$_{tỏa} = \frac {1{,}5.91{,}8}{3} = 45{,}9$ kJ.
\end{loigiaibth}
\def\writeANS{}
\begin{loigiaibth}{51}
  $n_{H_2} = 1$ mol ; $n_{O_2} = 1$ mol $\Rightarrow $ $H_2$ phản ứng hết, $O_2$ còn dư. \par Theo phương trình ta có cứ 2 mol $H_2$ phản ứng tỏa ra nhiệt lượng $572$ kJ \par $\Rightarrow $ 1 mol $H_2$ phản ứng tỏa ra nhiệt lượng Q = $\frac {1}{2}.\Delta _r H_{298}^\circ = 286$ kJ
\end{loigiaibth}
\def\writeANS{}
\begin{loigiaibth}{52}
  Ta có (2) = 2.(1) $\Rightarrow $ $\Delta _r H_{298}^\circ $ (2) = 2.280 = 560 kJ.
\end{loigiaibth}
\def\writeANS{}
\begin{loigiaibth}{53}
  Ta có : (2) = $\frac {-(1)}{2}$ $\Rightarrow $ $\Delta _r H_{298}^\circ (2) = \frac {-(-197)}{2} = 98{,}5$ kJ.
\end{loigiaibth}
\def\writeANS{}
\begin{loigiaibth}{54}
 Bao gồm: b, d. \par (a) Sai vì điều kiện chuẩn áp suất là 1 bar. \par (c) Sai vì phản ứng làm môi trường xung quanh nóng lên là phản ứng tỏa nhiệt.
\end{loigiaibth}
\def\writeANS{}
\begin{loigiaibth}{55}
 Bao gồm: b, c. \par (b) Sai vì phản ứng (2) không phải tạo thành từ các đơn chất. \par (c) Sai vì biến thiên enthalpy chuẩn của phản ứng giữa 1 mol $N_2$ với 1 mol $O_2$ tạo thành 2 mol NO là $\Delta _r H_{298(1)}^\circ $ kJ.
\end{loigiaibth}
\def\writeANS{}
\begin{loigiaibth}{56}
 Bao gồm: b, c, d, e. \par (a) Sai vì biến thiên enthalpy chuẩn của phản ứng là $-196{,}9$kJ
\end{loigiaibth}
\def\writeANS{}
\begin{loigiaibth}{57}
  Gọi số mol $C_3H_8$ và số mol $C_4H_{10}$ là $2a$, ta có: $44a + 58.2a = 12.1000 \Rightarrow a = 75$ mol \par Nhiệt đốt cháy 12 kg gas là Q = $75.2220 + 150.2874 = 597600$ (kJ) \par Số ngày sử dụng hết bình gas = $\frac {597600}{10000.\frac {100}{80}} = 47{,}808 \approx 48$ (ngày)
\end{loigiaibth}
\def\writeANS{}
\begin{loigiaibth}{58}
  Gọi số mol $CH_3OH$ và $C_2H_5OH$ trong 10g X lần lượt là a và b \par Ta có: $32a + 46b = 10$ (1) \par $716a + 1370b = 291{,}9$ (2) \par Giải (1) và (2), ta được: $a = 0{,}025; b = 0{,}2$ \par $\Rightarrow $ Khối lượng $CH_3OH$ là $32{,}0.025 = 0{,}8$ (g) \par $\Rightarrow $ Phần trăm tạp chất methanol trong X bằng $8\%$
\end{loigiaibth}
